\documentclass[11pt]{article}\usepackage{layout}\begin{document}
\begin{canvas}\Header{第二章 · 全球开源生态地图}{teal}\Footer{01}\pdfbookmark[1]{2.1 全球开源活跃地图}{page1}
\Box{15}{27}{180}{\fontsize{21}{27}\selectfont\color{ink}\sffamily \bfseries 2.1 全球开源活跃地图}
\Box{15}{40}{180}{\fontsize{9.7}{13.5}\selectfont\color{muted}\sffamily 活跃开发者指报告期内留下公开活动记录的账号；累计规模还包括当期未活跃账号。地图和排行采用 GitHub 口径，五域另用综合统计。}
\Box{15}{53}{180}{\includegraphics[width=180mm]{.build/assets/01-global-current.pdf}}
\Box{15}{199}{180}{\fontsize{7.8}{10.8}\selectfont\color{muted}\sffamily 图 2-1  全球参与全景与平台覆盖。来源：OpenDigger / OpenShare（GOSD 年度统计）；五域及国内平台补充由报告编制组提供。主统计期：2025.07—2026.06。}
\Box{15}{211}{87}{\fontsize{10.5}{14.2}\selectfont\color{teal}\sffamily \bfseries ① 多中心参与，构成全球合作基础}
\Box{15}{217}{87}{\fontsize{10.1}{15.6}\selectfont\color{ink}美国、印度、中国和巴西是本图规模较大的节点；欧洲由多个相邻国家共同构成参与网络。地图帮助发现合作基础，参与人数还不能直接代表技术质量、AI 研发实力或公共成效。}
\Box{108}{211}{87}{\fontsize{10.5}{14.2}\selectfont\color{teal}\sffamily \bfseries ② 排名之外，也要看见更多参与者}
\Box{108}{217}{87}{\fontsize{10.1}{15.6}\selectfont\color{ink}前 20 名把空间位置转为可比较的读数。除既有欧美节点外，印度、巴西、印度尼西亚等已形成可观参与规模；国际合作布局需要同时看见现有中心与成长中的社区。}
\Box{15}{245}{87}{\fontsize{10.5}{14.2}\selectfont\color{teal}\sffamily \bfseries ③ 用户基础扩大，持续协作更关键}
\Box{15}{251}{87}{\fontsize{10.1}{15.6}\selectfont\color{ink}2022—2026 年，累计规模增长 131.7\%，当期活跃增长 87.3\%。两者并不同步，生态评价需要继续观察贡献采纳和长期维护；这一差异不能解释为同一批用户的留存下降。}
\Box{108}{245}{87}{\fontsize{10.5}{14.2}\selectfont\color{teal}\sffamily \bfseries ④ 跨平台观察，补全中国参与图景}
\Box{108}{251}{87}{\fontsize{10.1}{15.6}\selectfont\color{ink}五域中国为 454.5 万，采用独立综合统计。GitHub 250.48 万加国内平台去重补充约 120 万，跨平台约 370 万；补充扩大了可见范围，不再叠加到五域读数。}
\end{canvas}
\begin{canvas}\Header{第二章 · 全球开源生态地图}{teal}\Footer{02}\pdfbookmark[1]{增长版图与 2030 情景}{page2}
\Box{15}{27}{180}{\fontsize{21}{27}\selectfont\color{ink}\sffamily \bfseries 增长版图与 2030 情景}
\Box{15}{40}{180}{\fontsize{9.7}{13.5}\selectfont\color{muted}\sffamily 实线为 2022—2026 报告年度；虚线为 2027—2030 条件情景。五域综合趋势与南方六国 GitHub 细分分别展示，覆盖不同。}
\Box{15}{53}{180}{\includegraphics[width=180mm]{.build/assets/02-trends-2030.pdf}}
\Box{15}{199}{180}{\fontsize{7.8}{10.8}\selectfont\color{muted}\sffamily 图 2-2  历史增长、2030 条件情景与南方六国细分。来源：OpenDigger / OpenShare、编制组五域序列及情景测算。情景不赋予概率；方法与完整读数见附录 1。}
\Box{15}{211}{87}{\fontsize{10.5}{14.2}\selectfont\color{teal}\sffamily \bfseries ① 增长与积累，形成不同的地域优势}
\Box{15}{217}{87}{\fontsize{10.1}{15.6}\selectfont\color{ink}2022—2026 年，南方六国综合组和印度的年均增速分别为 22.2\% 和 21.3\%，高于欧美比较组。若这些增长条件延续，其相对规模可能提高；欧洲仍保持较大的参与基础。}
\Box{108}{211}{87}{\fontsize{10.5}{14.2}\selectfont\color{teal}\sffamily \bfseries ② 把长期预测变成可更新的条件判断}
\Box{108}{217}{87}{\fontsize{10.1}{15.6}\selectfont\color{ink}三条小图分别追踪全球累计、全球活跃和中国 GitHub 活跃，避免把注册增长等同于持续参与。2030 范围来自“历史延续”和“逐步放缓”两种假设，每年应以同口径数据重算。}
\Box{15}{247}{180}{\fontsize{10.5}{14.2}\selectfont\color{teal}\sffamily \bfseries ③ 全球南方内部，合作需求同样具有差异}
\Box{15}{253}{180}{\fontsize{10.2}{16}\selectfont\color{ink}在六国 GitHub 细分中，巴西占活跃规模的 61.5\%；菲律宾和印度尼西亚的同比增速分别为 12.5\% 和 11.5\%，快于巴西的 7.4\%。AI 普惠合作需要同时支持已有节点与成长社区，由基金会、项目团队和本地伙伴共同提供本地化、导师与维护资源，并以实际采纳和持续贡献检验成效。}
\end{canvas}
\begin{canvas}\Header{第二章 · 全球开源生态地图}{teal}\Footer{03}\pdfbookmark[1]{中国区域参与结构}{page3}
\Box{15}{27}{180}{\fontsize{21}{27}\selectfont\color{ink}\sffamily \bfseries 中国区域参与结构}
\Box{15}{40}{180}{\fontsize{9.7}{13.5}\selectfont\color{muted}\sffamily 将现有参与基础、同比净增与增长速度联合观察。统一采用 GitHub 地理统计，国内平台补充未按省级分摊。}
\Box{15}{53}{180}{\includegraphics[width=180mm]{.build/assets/03-china-regions.pdf}}
\Box{15}{199}{180}{\fontsize{7.8}{10.8}\selectfont\color{muted}\sffamily 图 2-3  中国参与基础、同比增长与区域份额。来源：OpenDigger / OpenShare；省级值按平台地理估计阅读。同比为 2026 较 2025 报告年度；各份额独立取整。}
\Box{15}{211}{87}{\fontsize{10.5}{14.2}\selectfont\color{teal}\sffamily \bfseries ① 头部集聚与广泛参与并存}
\Box{15}{217}{87}{\fontsize{10.1}{15.6}\selectfont\color{ink}全国 GitHub 活跃规模约 250.48 万，同比净增 30.40 万。京沪粤占当期活跃的 52.6\%，形成主要参与基础；其余地区共同构成跨区域协作的重要支撑，不能仅以头部排行概括中国。}
\Box{108}{211}{87}{\fontsize{10.5}{14.2}\selectfont\color{teal}\sffamily \bfseries ② 现有规模与当期增量联合阅读}
\Box{108}{217}{87}{\fontsize{10.1}{15.6}\selectfont\color{ink}北京、上海、广东的净增分别约为 6.45 万、5.07 万和 4.39 万。前 20 表同时列出活跃与净增，帮助地方决策者区分参与基础和年度变化；净增并不等于人才迁入或首次参与。}
\Box{15}{245}{87}{\fontsize{10.5}{14.2}\selectfont\color{teal}\sffamily \bfseries ③ 浙江等地区展现更快增长}
\Box{15}{251}{87}{\fontsize{10.1}{15.6}\selectfont\color{ink}台湾、浙江的活跃规模同比增长 19.2\% 和 16.1\%，高于全国的 13.8\%；浙江同时拥有 23.29 万的活跃基础。省级差异为后续产业与人才观察提供线索，不能直接替代杭州等城市的证据。}
\Box{108}{245}{87}{\fontsize{10.5}{14.2}\selectfont\color{teal}\sffamily \bfseries ④ 增量结构提示区域合作机会}
\Box{108}{251}{87}{\fontsize{10.1}{15.6}\selectfont\color{ink}台湾与浙江占当期活跃的 18.9\%，占同比净增的 23.3\%，相差约 4.4 个百分点。地方支持应进一步连接产业项目、组织与维护团队，判断参与增长能否转化为持续供给和协作能力。}
\end{canvas}
\begin{canvas}\Header{附录 1 · 口径与来源}{teal}\Footer{A1}\pdfbookmark[1]{附录1 指标、口径与来源}{appendix1}
\Box{15}{28}{180}{\fontsize{20}{27}\selectfont\color{ink}\sffamily \bfseries 附录 1｜指标、口径与来源}
\Box{15}{43}{180}{\fontsize{10.3}{17}\selectfont\color{ink}2026 报告年度为 2025 年 7 月 1 日至 2026 年 6 月 30 日；2022—2025 年按同样的跨年窗口阅读。累计规模是期末基础，活跃规模覆盖整个报告年度。}
\Box{15}{66}{180}{\fontsize{12}{18}\selectfont\color{teal}\sffamily \bfseries 1  人数衡量什么}
\Box{15}{79}{29}{\fontsize{10.2}{16}\selectfont\color{teal}\sffamily \bfseries 累计规模}
\Box{46}{79}{149}{\fontsize{10}{16}\selectfont\color{ink}平台已有账号基础，包含当期没有公开活动的账号。不是职业开发者人口，也不代表其中每个人都参与 AI 研发。}
\Box{15}{100}{29}{\fontsize{10.2}{16}\selectfont\color{teal}\sffamily \bfseries 活跃开发者}
\Box{46}{100}{149}{\fontsize{10}{16}\selectfont\color{ink}报告期内出现在公开活动记录中的账号，活动包括收藏项目、复制仓库、提交修改或讨论等；不能与“代码贡献者”直接互换。}
\Box{15}{124}{29}{\fontsize{10.2}{16}\selectfont\color{teal}\sffamily \bfseries 同比与净增}
\Box{46}{124}{149}{\fontsize{10}{16}\selectfont\color{ink}净增 = 本期活跃 − 上期活跃；同比增速 = 净增 ÷ 上期活跃。它们不直接测量首次参与、人才迁移或同一批人的留存。}
\Box{15}{145}{29}{\fontsize{10.2}{16}\selectfont\color{teal}\sffamily \bfseries 地理分布}
\Box{46}{145}{149}{\fontsize{10}{16}\selectfont\color{ink}地区值包含按可识别所在地进行的估计，故原始读数可能有小数；表示平台可观察的地域参与，不表示国籍、职业或雇主归属。}
\Box{15}{170}{180}{\fontsize{12}{18}\selectfont\color{teal}\sffamily \bfseries 2  三套覆盖分别阅读}
\Box{15}{183}{180}{\fontsize{9.7}{15.4}\selectfont\color{ink}国家地图、前 20、全球与中国三条趋势及省级图：采用 OpenShare 的 GOSD 年度 GitHub 数据；南方六国细分与之同源。}
\Box{15}{202}{180}{\fontsize{9.7}{15.4}\selectfont\color{ink}五域：欧洲、美国、中国、印度、南方六国采用报告编制组的独立综合序列，不是全球完整分区。五域与单国平台表的覆盖不同，不混排、不相加。}
\Box{15}{226}{180}{\fontsize{9.7}{15.4}\selectfont\color{ink}中国跨平台：GitHub 250.48 万，加 Gitee、AtomGit 等国内平台经问卷及平台数据去重补充约 120 万，合计约 370 万。该补充已确认；不回填历史、不分摊到省，也不再次叠加到五域。}
\Box{15}{252}{180}{\fontsize{8.7}{13}\selectfont\color{muted}来源：OpenDigger / OpenShare、报告编制组。公开数据读取于 2026-11-09；其元数据更新时间为 2026-10-08。地图几何沿用 2026 版，台湾、香港、澳门纳入中国。}
\Box{15}{272}{180}{\sffamily\fontsize{8.4}{11}\selectfont\href{https://open-share.com/insight/open-world/gosd}{OpenShare GOSD 入口}\quad\href{https://selfoss.open-digger.cn/openshare/overview_half_year/meta.json}{统计期元数据}\quad\href{https://selfoss.open-digger.cn/openshare/overview_half_year/developers.json}{全球及国家数据}\quad\href{https://selfoss.open-digger.cn/openshare/overview_half_year/CN/developers.json}{中国省级数据}}
\end{canvas}
\begin{canvas}\Header{附录 1 · 历史与情景}{teal}\Footer{A2}\pdfbookmark[2]{历史序列、情景测算与评审记录}{appendix2}
\Box{15}{28}{180}{\fontsize{20}{27}\selectfont\color{ink}\sffamily \bfseries 历史序列与 2030 条件情景}
\Box{15}{44}{180}{\fontsize{9.9}{16}\selectfont\color{ink}南方六国指巴西、印度尼西亚、越南、菲律宾、埃及、南非，不代表全球南方整体；印度单列。上方五域综合数不等同于六国 GitHub 细分之和。}
\Box{15}{65}{180}{\fontsize{11.8}{17}\selectfont\color{teal}\sffamily \bfseries 3  历史读数｜2022—2026 报告年度}
\Box{15}{77}{180}{\sffamily\fontsize{9.3}{14}\selectfont\renewcommand{\arraystretch}{1.12}\setlength{\tabcolsep}{5.1pt}\begin{tabularx}{180mm}{@{}Xrrrrr@{}}\toprule 对象／单位 & 2022 & 2023 & 2024 & 2025 & 2026\\\midrule
欧洲／万&370.35&414.79&464.56&520.48&675.39\\
美国／万&252.43&282.72&316.65&354.60&465.15\\
中国／万&237.79&266.80&299.60&320.00&454.53\\
印度／万&125.90&149.80&178.30&212.20&272.36\\
南方六国／万&134.60&161.50&193.80&232.62&300.46\\
全球累计／亿&0.84&1.05&1.29&1.57&1.94\\
全球活跃／万&1,543.27&1,794.31&2,009.11&2,644.05&2,890.37\\
中国 GitHub／万&129.47&145.71&162.47&220.08&250.48\\
\bottomrule\end{tabularx}}
\Box{15}{135}{180}{\fontsize{11.8}{17}\selectfont\color{teal}\sffamily \bfseries 4  2030 情景读数｜不是确定结果或置信区间}
\Box{15}{147}{180}{\sffamily\fontsize{9.3}{14}\selectfont\renewcommand{\arraystretch}{1.12}\setlength{\tabcolsep}{7pt}\begin{tabularx}{180mm}{@{}Xrrr@{}}\toprule 对象／单位 & 历史年均增速 & 2030 放缓 & 2030 延续\\\midrule
欧洲／万&16.21\%&972.53&1,231.67\\
美国／万&16.51\%&674.16&857.15\\
中国／万&17.58\%&674.01&868.83\\
印度／万&21.28\%&436.56&589.19\\
南方六国／万&22.23\%&491.27&670.69\\
全球累计／亿&23.38\%&3.25&4.49\\
全球活跃／万&16.98\%&4,231.74&5,413.33\\
中国 GitHub／万&17.94\%&374.24&484.59\\
\bottomrule\end{tabularx}}
\Box{15}{205}{180}{\fontsize{9.6}{15.1}\selectfont\color{ink}以 2026 年为起点，历史年均增速 g =（2026 人数 ÷ 2022 人数）的四次方根 − 1。延续情景每年按 g 增长；放缓情景在第 h 年采用 g × 0.8 的 h 次方，逐年复利计算。测算使用未舍入值。}
\Box{15}{228}{180}{\fontsize{9.6}{15.1}\selectfont\color{ink}0.8 是预设的敏感性参数，未拟合概率，也未单独估计 AI 或政策的因果效应。每年补入新观测后重算；发布后 6—12 个月优先跟踪贡献采纳、持续参与和维护职责。}
\Box{15}{251}{180}{\fontsize{8.5}{12.4}\selectfont\color{muted}评审记录｜v0.3，2026-11-09：三页正文全部改为 A4 竖版；恢复中国五域 454.534 万，纳入已确认的跨平台约 370 万；补齐三条趋势、南方六国和中国增量分析。}
\Box{15}{268}{180}{\fontsize{8.3}{12}\selectfont\color{muted}待核查：五域欧洲成员清单及历史覆盖执行记录；地图出版规范。待评审：竖版图文密度与情景假设。上述事项不改变本轮已确认的中国补充数。}
\end{canvas}\end{document}